% Section 9.1, from lectures/L09/appendix/a_structs_as_drivers.md. The gpio_driver_demo's GPIO
% driver and main.c are typeset straight from their files in §9.1.5; its EEPROM, serial and
% watchdog drivers are the ones printed in Chapters 6 to 8.

\section{Structs as Drivers}
\label{c9:sec:structs}\label{c9:app:a}

\subsection{Why bundle state}
\label{c9:a:1}
\begin{itemize}
  \item The drivers covered in Chapters~\ref{c4:ch} and \ref{c6:ch}--\ref{c8:ch} --- the ADC,
    UART, EEPROM, and the watchdog --- did not need any structs:
  \begin{itemize}
    \item The chip has exactly one ADC (\code{adc_read()} already takes a channel number as a
      parameter, so it was never hardcoded to a single pin the way the examples below are).
    \item It has exactly one UART.
    \item It has exactly one EEPROM.
    \item It has exactly one watchdog.
    \item Because there is only one instance of each peripheral, plain functions operating on
      implicit global state were the simplest correct solution.
  \end{itemize}
  \item The same single-instance pattern shows up again with peripherals the chip actually has
    many of, where it isn't the right solution:
  \begin{itemize}
    \item \code{LED_ON} assumes one specific LED, even though the chip has many GPIO pins.
    \item The configuration is effectively hardcoded into macros, functions, and global state.
  \end{itemize}
  \item This pattern works well as long as the program needs only one instance of each thing.
  \item The pattern starts to break down when the program needs several instances:
  \begin{itemize}
    \item \Secref{c4:app:b} contains an exercise with five LEDs, Exercise~\ref{c4:ex:7}.
    \item Without a better abstraction, adding another LED means copying an entire family of
      macros or functions and renaming them.
    \item That quickly becomes repetitive, difficult to maintain, and easy to get wrong.
  \end{itemize}
  \item C already provides the tool needed to solve this problem: the \code{struct}:
  \begin{itemize}
    \item Store one peripheral instance's configuration and state in a \code{struct}.
    \item Create one struct value for each LED (or other multi-instance peripheral) instance.
    \item Write ordinary functions that receive a pointer to the relevant struct.
    \item The same functions can then operate on several different instances.
  \end{itemize}
  \item By convention, the pointer parameter is called \code{self}:
  \begin{itemize}
    \item \code{self} points to the particular peripheral instance that the function should
      operate on.
    \item It plays a role similar to \code{this} in languages with built-in object support, such
      as C++, C\#, and Java.
    \item Unlike \code{this}, however, \code{self} is just an ordinary function parameter written
      explicitly in the code.
  \end{itemize}
  \item However, C has no built-in support for object-oriented programming:
  \begin{itemize}
    \item There is no hidden behaviour.
    \item There are no constructors.
    \item There is no inheritance.
    \item There is no automatic method dispatch.
    \item There is only a \code{struct} containing data and a set of ordinary functions that
      accept a pointer to that data.
  \end{itemize}
  \item This chapter keeps the design deliberately simple:
  \begin{itemize}
    \item The struct's fields remain directly accessible.
    \item Functions are called normally.
    \item No function pointers are required.
  \end{itemize}
  \item The next chapter, Chapter~\ref{c10:ch}, explores how much further this approach can be
    taken:
  \begin{itemize}
    \item Hiding a struct's internal fields.
    \item Calling operations through function pointers.
    \item Creating a more object-like interface while still using C.
  \end{itemize}
\end{itemize}

\subsection{The \code{self} convention: a \code{GPIO} driver}
\label{c9:a:2}
A minimal driver is a struct holding whatever a function needs to know to act on one specific
instance: here, pointers to the GPIO's data direction (\code{DDRx}), output (\code{PORTx}), and
input (\code{PINx}) registers, plus the bit position of its pin within them.

See \secref{c1:a:3} for a refresher on structs and \code{typedef}; \code{gpio_t} follows the same
anonymous-struct-plus-\code{_t}-suffix convention introduced there.

\begin{ccode}
/**
 * @brief GPIO driver structure.
 */
typedef struct
{
    /** Pointer to the data direction register associated with the GPIO. */
    volatile uint8_t* ddrx;

    /** Pointer to the output register associated with the GPIO. */
    volatile uint8_t* portx;

    /** Pointer to the input register associated with the GPIO. */
    volatile uint8_t* pinx;

    /** Bit position of the GPIO pin within the associated port. */
    const uint8_t pin;
} gpio_t;

/**
 * @brief Enumeration of GPIO modes.
 */
typedef enum
{
    GPIO_MODE_INPUT,        ///< Standard GPIO input.
    GPIO_MODE_INPUT_PULLUP, ///< GPIO input with its internal pull-up enabled.
    GPIO_MODE_OUTPUT,       ///< GPIO output.
} gpio_mode_t;

/**
 * @brief Initialize GPIO.
 *
 * @param[in] self GPIO instance.
 * @param[in] mode GPIO mode.
 */
void gpio_init(gpio_t* self, const gpio_mode_t mode)
{
    if (NULL == self) { return; }

    switch (mode)
    {
        case GPIO_MODE_INPUT_PULLUP:
            *self->ddrx &= ~(1U << self->pin);
            *self->portx |= (1U << self->pin);
            break;
        case GPIO_MODE_OUTPUT:
            *self->ddrx |= (1U << self->pin);
            *self->portx &= ~(1U << self->pin);
            break;
        default:
            *self->ddrx &= ~(1U << self->pin);
            *self->portx &= ~(1U << self->pin);
            break;
    }
}

/**
 * @brief Set state of the given GPIO.
 *
 * @param[in] self GPIO instance.
 * @param[in] state GPIO state (true = high, false = low).
 *
 * @note This operation is only supported for GPIOs configured as outputs.
 */
void gpio_write(gpio_t* self, const bool state)
{
    if (NULL == self) { return; }
    const bool output = (*self->ddrx & (1U << self->pin));
    if (!output) { return; }

    if (state) { *self->portx |= (1U << self->pin); }
    else { *self->portx &= ~(1U << self->pin); }
}

/**
 * @brief Read state of the given GPIO.
 *
 * @param[in] self GPIO instance.
 *
 * @return GPIO state (true = high, false = low).
 */
bool gpio_read(const gpio_t* self)
{
    return (NULL != self) ? (*self->pinx & (1U << self->pin)) : false;
}

/**
 * @brief Toggle state of the given GPIO.
 *
 * @param[in] self GPIO instance.
 */
void gpio_toggle(gpio_t* self)
{
    if (NULL == self) { return; }
    const bool output = (*self->ddrx & (1U << self->pin));
    if (!output) { return; }
    *self->pinx = (1U << self->pin);
}
\end{ccode}

\noindent A GPIO output instance can then be initialized as shown below:

\begin{ccode}
#define LED1_PIN 1U // D9 -> PORTB1.

/** Initialize LED on D9 (PORTB1). */
gpio_t led1 = {&DDRB, &PORTB, &PINB, LED1_PIN};

gpio_init(&led1, GPIO_MODE_OUTPUT);
\end{ccode}

\noindent Most of these functions take \code{gpio_t*} rather than \code{gpio_t}: passing a pointer
avoids copying the struct on every call, and they need to write through \code{self} to the
hardware registers it points at. \code{gpio_read()} takes \code{const gpio_t*} instead, since it
only ever reads \code{self}, never writes through it.

\subsection{Splitting the driver into \code{gpio.h} and \code{gpio.c}}
\label{c9:a:3}
So far \secref{c9:a:2}'s \code{gpio_t}, \code{gpio_mode_t}, and every \code{gpio_*()} function
have lived together in one code block. In a real project a driver is split across a header and a
source file, the same split \secref{c1:a:6} introduced for \file{sensor.h}/\file{sensor.c},
include guard and all:
\begin{itemize}
  \item The header (\file{gpio.h}) is the public interface: the \code{gpio_t} and
    \code{gpio_mode_t} type definitions, plus a documented prototype for every function callers
    are allowed to use.
  \item The source file (\file{gpio.c}) \code{#include}s its own header, after the standard and
    AVR headers as everywhere in this book, then defines those same functions. Callers only ever
    see the header; the definitions in \file{gpio.c} are what actually gets compiled.
\end{itemize}
One thing is new compared to \file{sensor.h}: that header only ever declared a function, callers
never touched anything but the value it returned. \file{gpio.h} also declares \code{gpio_t}
itself, because a caller needs to create and initialize their own \code{gpio_t} instances directly
(as in \secref{c9:a:2}'s \code{led1} example), not just receive values back from an opaque handle.
That's also why \code{gpio_t}'s fields stay public here instead of hidden behind a pointer to an
incomplete type; Chapter~\ref{c10:ch} revisits that trade-off.

Here's \file{driver/gpio.h}:

\begin{ccode}
/**
 * @file GPIO driver.
 */
#ifndef DRIVER_GPIO_H_
#define DRIVER_GPIO_H_

#include <stdbool.h>
#include <stdint.h>

/**
 * @brief GPIO driver structure.
 */
typedef struct
{
    /** Pointer to the data direction register associated with the GPIO. */
    volatile uint8_t* ddrx;

    /** Pointer to the output register associated with the GPIO. */
    volatile uint8_t* portx;

    /** Pointer to the input register associated with the GPIO. */
    volatile uint8_t* pinx;

    /** Bit position of the GPIO pin within the associated port. */
    const uint8_t pin;
} gpio_t;

/**
 * @brief Enumeration of GPIO modes.
 */
typedef enum
{
    GPIO_MODE_INPUT,        ///< Standard GPIO input.
    GPIO_MODE_INPUT_PULLUP, ///< GPIO input with its internal pull-up enabled.
    GPIO_MODE_OUTPUT,       ///< GPIO output.
} gpio_mode_t;

/**
 * @brief Initialize GPIO.
 *
 * @param[in] self GPIO instance.
 * @param[in] mode GPIO mode.
 */
void gpio_init(gpio_t* self, gpio_mode_t mode);

/**
 * @brief Set state of the given GPIO.
 *
 * @param[in] self GPIO instance.
 * @param[in] state GPIO state (true = high, false = low).
 *
 * @note This operation is only supported for GPIOs configured as outputs.
 */
void gpio_write(gpio_t* self, bool state);

/**
 * @brief Read state of the given GPIO.
 *
 * @param[in] self GPIO instance.
 *
 * @return GPIO state (true = high, false = low).
 */
bool gpio_read(const gpio_t* self);

/**
 * @brief Toggle state of the given GPIO.
 *
 * @param[in] self GPIO instance.
 */
void gpio_toggle(gpio_t* self);

#endif /** DRIVER_GPIO_H_ */
\end{ccode}

\noindent Here's \file{driver/gpio.c}:

\begin{ccode}
/**
 * @file GPIO driver implementation details.
 */
#include <stdbool.h>
#include <stddef.h>
#include <stdint.h>

#include <avr/io.h>

#include "driver/gpio.h"

// -----------------------------------------------------------------------------
void gpio_init(gpio_t* self, const gpio_mode_t mode)
{
    // Check the GPIO instance, terminate if invalid.
    if (NULL == self) { return; }

    // Configure GPIO mode as specified.
    switch (mode)
    {
        case GPIO_MODE_INPUT_PULLUP:
            // Configure GPIO as input with its internal pull-up enabled.
            *self->ddrx &= ~(1U << self->pin);
            *self->portx |= (1U << self->pin);
            break;
        case GPIO_MODE_OUTPUT:
            // Configure GPIO as output.
            *self->ddrx |= (1U << self->pin);
            *self->portx &= ~(1U << self->pin);
            break;
        default:
            // Configure GPIO as standard input.
            *self->ddrx &= ~(1U << self->pin);
            *self->portx &= ~(1U << self->pin);
            break;
    }
}

// -----------------------------------------------------------------------------
void gpio_write(gpio_t* self, const bool state)
{
    // Check the GPIO instance, terminate if invalid.
    if (NULL == self) { return; }

    // Check GPIO direction, terminate if not output.
    const bool output = (*self->ddrx & (1U << self->pin));
    if (!output) { return; }

    // Set GPIO state as specified.
    if (state) { *self->portx |= (1U << self->pin); }
    else { *self->portx &= ~(1U << self->pin); }
}

// -----------------------------------------------------------------------------
bool gpio_read(const gpio_t* self)
{
    // Return GPIO state if GPIO instance is valid, else false.
    return (NULL != self) ? (*self->pinx & (1U << self->pin)) : false;
}

// -----------------------------------------------------------------------------
void gpio_toggle(gpio_t* self)
{
    // Check the GPIO instance, terminate if invalid.
    if (NULL == self) { return; }

    // Check GPIO direction, terminate if not output.
    const bool output = (*self->ddrx & (1U << self->pin));
    if (!output) { return; }

    // Toggle GPIO state.
    *self->pinx = (1U << self->pin);
}
\end{ccode}

\noindent Notice \code{gpio_init()}'s \code{mode} parameter (and \code{gpio_write()}'s
\code{state}) drop the \code{const} in the header but keep it in the source file. A top-level
\code{const} on a parameter only constrains what the function body can do with its own local copy;
it isn't part of the function's type and can't affect how callers call it, so by convention it's
omitted from the prototype and kept in the definition, where it actually matters.

\subsection{Worked example: an array of LEDs}
\label{c9:a:4}
The whole point of a driver struct is that \code{gpio_write()}/\code{gpio_read()} don't care which
GPIO they're called on. An array of \code{gpio_t} instances plus a loop replaces a whole family of
per-LED macros:

\begin{ccode}
#define LED_COUNT 5U

// Create five LED instances.
static gpio_t leds[LED_COUNT] =
{
    { &DDRD, &PORTD, &PIND, 6U}, // D6  -> PORTD6.
    { &DDRD, &PORTD, &PIND, 7U}, // D7  -> PORTD7.
    { &DDRB, &PORTB, &PINB, 0U}, // D8  -> PORTB0.
    { &DDRB, &PORTB, &PINB, 1U}, // D9  -> PORTB1.
    { &DDRB, &PORTB, &PINB, 2U}, // D10 -> PORTB2.
};

/**
 * @brief Set up system.
 */
static void setup(void)
{
    for (uint8_t i = 0U; i < LED_COUNT; i++)
    {
        gpio_init(&leds[i], GPIO_MODE_OUTPUT);
    }
}
\end{ccode}

\noindent Compare this to \secref{c4:app:b}'s five-LED exercise, Exercise~\ref{c4:ex:7}, which
needed a hand-written \code{LED1_ON} through \code{LED5_ON} (and \code{_OFF}), one pair of macros
per pin. Here, adding a sixth LED on the same port means adding one line to the \code{leds[]}
array, no new functions.

\subsection{Worked example: \code{gpio_t} wired into Chapters 6--8's demo}
\label{c9:a:5}
Combines Chapter~\ref{c6:ch}'s serial driver, Chapter~\ref{c7:ch}'s EEPROM driver, and
Chapter~\ref{c8:ch}'s watchdog driver with this section's \code{gpio_t}: LED1 (D9) persists its
state across resets via EEPROM, BTN1 (D13) toggles LED1 through a debounced pin-change interrupt
exactly like \secref{c8:a:5}, and the watchdog resets the chip if the main loop ever stalls. Every
register access that went through Chapter~\ref{c8:ch}'s \code{LED1_TOGGLE}/\code{LED1_ENABLED}/%
\code{BTN1_PRESSED} macros now goes through \code{gpio_toggle()}/\code{gpio_read()} instead.

The circuit is Chapter~\ref{c6:ch}'s, in Figure~\ref{c6:fig:serial}.

\begin{console}
gpio_driver_demo/
├── Makefile
├── main.c
├── include/
│   └── driver/
│       ├── eeprom.h
│       ├── gpio.h
│       ├── serial.h
│       └── watchdog.h
└── source/
    └── driver/
        ├── eeprom.c
        ├── gpio.c
        ├── serial.c
        └── watchdog.c
\end{console}

\begin{itemize}
  \item The serial, EEPROM, and watchdog drivers were introduced in
    Chapters~\ref{c6:ch}--\ref{c8:ch}.
  \item \file{driver/gpio.h} and \file{driver/gpio.c}, both printed in full below, extend the
    \code{gpio_t} type introduced in \secref{c9:a:2} and \secref{c9:a:3}.
  \item The extended driver adds one capability not covered in those sections:
  \begin{itemize}
    \item \code{gpio_enable_pci()} enables or disables the pin-change interrupt for a GPIO.
    \item It does this by writing the same \code{PCICR} and \code{PCMSKx} bits that were
      configured manually in \secref{c8:a:5}.
  \end{itemize}
  \item Supporting \code{gpio_enable_pci()} requires two additional fields in \code{gpio_t},
    beyond \code{ddrx}, \code{portx}, \code{pinx}, and \code{pin}:
  \begin{itemize}
    \item A pointer to the port's \code{PCMSKx} register.
    \item The bit position of the port's \code{PCIEx} bit in \code{PCICR}.
  \end{itemize}
  \item Every \code{gpio_t} instance contains these two additional fields:
  \begin{itemize}
    \item This is true even when that particular GPIO never uses \code{gpio_enable_pci()}.
    \item For example, LED1 still carries the pin-change-interrupt configuration fields.
  \end{itemize}
\end{itemize}

\subsubsection{The driver in full}
\cfile{lectures/L09/gpio_driver_demo/include/driver/gpio.h}
\needspace{8\baselineskip}
\cfile{lectures/L09/gpio_driver_demo/source/driver/gpio.c}

\subsubsection{\code{main.c} in full}
\cfile{lectures/L09/gpio_driver_demo/main.c}

\noindent\code{led1_event}, the debounce timer, and the EEPROM-persistence logic are unchanged from
\secref{c8:a:5}; only the GPIO handling changed, from bespoke macros tied to \code{PORTB}
specifically to a driver that works for however many \code{gpio_t} instances a program declares,
on whichever pins and ports they happen to be wired to.

\subsubsection{Building and running}
Build and flash the same way as \secref{c8:a:5}: on Linux, run \code{make} in the demo's
directory, \file{lectures/L09/gpio_driver_demo} (\code{make flash} to program the board); on
Windows via Microchip Studio, recreate this folder structure inside the project and add
\file{include} as an include directory.

The demo's \file{Makefile} is the one printed in \secref{c6:sec:build}, with
\file{source/driver/eeprom.c}, \file{source/driver/gpio.c} and \file{source/driver/watchdog.c}
added to its \code{SRC_FILES}. Its serial, EEPROM and watchdog drivers are byte for byte the ones
printed in \secref{c6:a:6}, \secref{c7:a:4} and \secref{c8:a:5}.
